\begin{pastexam}[label={px:ccu-108-swa-04}]{中正資工所甲組 108 學年度 第 5 題}
（填表）

(2\%) Assume that we have a 17-bucket hash table with one slot per bucket. The hash function for
each key using the division hash function, \textit{key}\%17. Sequentially insert the records with
keys are 6, 12, 34, 29, 28, 11, 23, 7, 0, and 33. Please complete the hash table at the end with
Linear Probing to handle overflows.

\begin{center}\small
\begin{tabular}{|*{17}{c|}}
\multicolumn{1}{c}{0} & \multicolumn{1}{c}{1} & \multicolumn{1}{c}{2} & \multicolumn{1}{c}{3} & \multicolumn{1}{c}{4} & \multicolumn{1}{c}{5} & \multicolumn{1}{c}{6} & \multicolumn{1}{c}{7} & \multicolumn{1}{c}{8} & \multicolumn{1}{c}{9} & \multicolumn{1}{c}{10} & \multicolumn{1}{c}{11} & \multicolumn{1}{c}{12} & \multicolumn{1}{c}{13} & \multicolumn{1}{c}{14} & \multicolumn{1}{c}{15} & \multicolumn{1}{c}{16}\\\hline
\phantom{00} & \phantom{00} & \phantom{00} & \phantom{00} & \phantom{00} & \phantom{00} & \phantom{00} & \phantom{00} & \phantom{00} & \phantom{00} & \phantom{00} & \phantom{00} & \phantom{00} & \phantom{00} & \phantom{00} & \phantom{00} & \phantom{00}\\\hline
\end{tabular}
\end{center}
\end{pastexam}

\begin{pastsoln}
線性探測（資結書第~12.3~節）：先算 $h(k) = k \bmod 17$，那一格有人就往後一格一格找空位，
超過第 $16$ 格繞回第 $0$ 格。

\begin{center}
\begin{tabular}{rrl}
\toprule
鍵 & $k \bmod 17$ & 探測的格子 $\to$ 放在 \\
\midrule
6 & 6 & $6 \to 6$ \\
12 & 12 & $12 \to 12$ \\
34 & $34 - 34 = 0$ & $0 \to 0$ \\
29 & $29 - 17 = 12$ & $12, 13 \to 13$ \\
28 & $28 - 17 = 11$ & $11 \to 11$ \\
11 & 11 & $11, 12, 13, 14 \to 14$ \\
23 & $23 - 17 = 6$ & $6, 7 \to 7$ \\
7 & 7 & $7, 8 \to 8$ \\
0 & 0 & $0, 1 \to 1$ \\
33 & $33 - 17 = 16$ & $16 \to 16$ \\
\bottomrule
\end{tabular}
\end{center}

幾個要注意的地方：$29$ 和 $12$ 撞在第 $12$ 格，往後放到第 $13$ 格；$11$ 的家是第 $11$ 格，
但前面的 $28$ 已經放在那裡，第 $12$、$13$ 格也被 $12$、$29$ 佔了，一路探到第 $14$ 格；
$7$ 的家第 $7$ 格被先來的 $23$ 佔走，放到第 $8$ 格。 最後的表：

\begin{center}\small
\begin{tabular}{|*{17}{c|}}
\multicolumn{1}{c}{0} & \multicolumn{1}{c}{1} & \multicolumn{1}{c}{2} & \multicolumn{1}{c}{3} & \multicolumn{1}{c}{4} & \multicolumn{1}{c}{5} & \multicolumn{1}{c}{6} & \multicolumn{1}{c}{7} & \multicolumn{1}{c}{8} & \multicolumn{1}{c}{9} & \multicolumn{1}{c}{10} & \multicolumn{1}{c}{11} & \multicolumn{1}{c}{12} & \multicolumn{1}{c}{13} & \multicolumn{1}{c}{14} & \multicolumn{1}{c}{15} & \multicolumn{1}{c}{16}\\\hline
34 & 0 & & & & & 6 & 23 & 7 & & & 28 & 12 & 29 & 11 & & 33\\\hline
\end{tabular}
\end{center}
\end{pastsoln}
